\documentclass[assignment.tex]{subfiles}
\begin{document}

\chapter{Assignment 3: RAG-Based Legal Contract Chatbot}
\label{sec:assignment3}

The third stage integrates all upstream outputs — structured clauses, NER entities,
SRL roles, and intent labels — into an end-to-end Retrieval-Augmented Generation
(RAG) chatbot.  Users can ask natural-language questions about a legal contract and
receive answers grounded in specific, cited clauses.

\section{System Overview}
\label{subsec:rag-overview}

The system is composed of three principal modules that form a linear pipeline
(Table~\ref{tab:rag-components}): a vector store for dense retrieval, an LCEL
RAG pipeline for answer generation, and a Streamlit front-end for user interaction.

\begin{table}[H]
    \centering
    \caption{RAG system components and their roles.}
    \label{tab:rag-components}
    \begin{tabular}{@{}lll@{}}
        \toprule
        \textbf{Module} & \textbf{Technology} & \textbf{Role} \\
        \midrule
        Embedding model  & \texttt{all-MiniLM-L6-v2}          & Dense clause representation \\
        Vector database  & ChromaDB (cosine similarity)        & Semantic clause retrieval \\
        LLM              & Gemini \texttt{gemini-1.5-flash}    & Grounded answer generation \\
        Orchestration    & LangChain LCEL                      & Prompt + chain management \\
        User interface   & Streamlit                           & Interactive web application \\
        \bottomrule
    \end{tabular}
\end{table}

\subsection{Vector Store and Indexing}
\label{subsec:vector-store}

\subsection{Embedding Model}

Clause embeddings are produced by
\texttt{sentence-transformers/all-MiniLM-L6-v2}, a lightweight 384-dimension
bi-encoder optimised for semantic similarity tasks.  Its small size makes it
suitable for real-time encoding of query strings at inference time.

\subsection{ChromaDB Configuration}

The vector store is implemented using ChromaDB's \texttt{PersistentClient},
enabling the index to survive application restarts without re-processing the
contract.  Cosine similarity is used as the distance metric:

\begin{lstlisting}[language=Python,
                   caption={\texttt{vector\_store.py} — ChromaDB collection
                   initialisation.},
                   label={lst:chroma-init}]
import chromadb
from chromadb.config import Settings
from sentence_transformers import SentenceTransformer

encoder = SentenceTransformer("sentence-transformers/all-MiniLM-L6-v2")

client = chromadb.PersistentClient(path="./chroma_db")
collection = client.get_or_create_collection(
    name="legal_clauses",
    metadata={"hnsw:space": "cosine"},
)
\end{lstlisting}

\subsection{Metadata Schema}

Each document stored in ChromaDB carries the following metadata fields derived
from the upstream pipeline outputs:

\begin{itemize}[leftmargin=*, itemsep=2pt]
    \item \texttt{ner\_entities} — JSON-serialised list of entity spans
          (text, label, start, end).
    \item \texttt{srl\_predicate} — main verb of the clause.
    \item \texttt{srl\_roles} — JSON-serialised role dictionary
          (Agent, Theme, Recipient, Time, Condition).
    \item \texttt{intent\_tfidf} — intent label from the TF-IDF classifier.
    \item \texttt{intent\_bert} — intent label from the BERT classifier.
\end{itemize}

Storing this metadata alongside the clause text enables the front-end to present
NER and SRL information for each retrieved source without additional inference.

\section{RAG Pipeline}
\label{subsec:rag-pipeline}

\subsection{LangChain LCEL Chain}

The retrieval and generation steps are orchestrated via LangChain's composable
LCEL (LangChain Expression Language) syntax:

\begin{lstlisting}[language=Python,
                   caption={\texttt{rag\_pipeline.py} — LCEL chain definition.},
                   label={lst:lcel-chain}]
from langchain_core.prompts import ChatPromptTemplate
from langchain_google_genai import ChatGoogleGenerativeAI
from langchain_core.output_parsers import StrOutputParser

llm = ChatGoogleGenerativeAI(model="gemini-1.5-flash", temperature=0.2)

chain = ChatPromptTemplate.from_messages([
    ("system", SYSTEM_PROMPT),
    ("human",  "{question}"),
]) | llm | StrOutputParser()
\end{lstlisting}

\subsection{Anti-Hallucination System Prompt}

A carefully designed system prompt constrains the model to stay within the
retrieved context.  The key directives are:

\begin{enumerate}[leftmargin=*, itemsep=3pt]
    \item \textbf{Citation requirement.}  Every factual claim must be supported by
          an explicit clause number (e.g.\ \textit{``As stated in Clause 7,
          \ldots''}).
    \item \textbf{Refusal policy.}  If the answer to the user's question cannot be
          found in the retrieved clauses, the model must explicitly decline to
          answer rather than confabulate.
    \item \textbf{Scope restriction.}  The model must answer only from the provided
          contract context and must not import external legal knowledge as a
          primary source.
\end{enumerate}

\begin{lstlisting}[language=Python,
                   caption={System prompt enforcing citation and refusal.},
                   label={lst:sys-prompt}]
SYSTEM_PROMPT = """You are a legal contract assistant.
Answer questions ONLY based on the contract clauses provided below.

Rules:
- Cite clause numbers for every claim (e.g. "Clause 3 states that...").
- If the answer is NOT found in the provided clauses, respond:
  "I cannot find information about this in the contract."
- Do NOT use external legal knowledge as a primary answer source.

Contract clauses:
{context}
"""
\end{lstlisting}

\subsection{Multi-Turn Conversation}

The \texttt{chat()} method maintains a list of \texttt{HumanMessage} and
\texttt{AIMessage} objects to support multi-turn dialogue.  The conversation
history is prepended to each new request, allowing the model to resolve
coreferences across turns (e.g.\ \textit{``Who does \textit{they} refer to in the
previous answer?''}).

\subsection{Top-$k$ Retrieval}

The number of retrieved clauses is configurable via the \texttt{top\_k} parameter
(default: 5).  At query time, the user's question is embedded with the same
\texttt{all-MiniLM-L6-v2} encoder, and the $k$ clauses with the highest cosine
similarity scores are returned from ChromaDB and injected into the system prompt
as the context window.

\section{Streamlit User Interface}
\label{subsec:streamlit}

\subsection{Interface Layout}

The web application (\texttt{app.py}) is built with Streamlit and organised into
two panels:

\begin{itemize}[leftmargin=*, itemsep=3pt]
    \item \textbf{Sidebar.}
    \begin{itemize}
        \item Top-$k$ slider (range 1--10) to control retrieval breadth.
        \item \textit{Build Index} button: triggers contract parsing and ChromaDB
              indexing.
        \item \textit{Clear Index} button: deletes the persisted collection to
              allow re-indexing with a different document.
    \end{itemize}
    \item \textbf{Main panel.}
    \begin{itemize}
        \item Chat input box for natural-language questions.
        \item Conversation history rendered as alternating user/assistant message
              bubbles.
        \item Colour-coded intent badges next to each retrieved source clause
              (Obligation = blue, Prohibition = red, Right = green,
              Termination Condition = orange).
        \item Expandable \textit{Sources} panel for each assistant response,
              revealing the retrieved clause text along with its NER entities and
              SRL roles extracted in the upstream pipeline.
    \end{itemize}
\end{itemize}

\subsection{Index Building Workflow}

When the user clicks \textit{Build Index}, the application executes the following
steps:

\begin{enumerate}[leftmargin=*, itemsep=3pt]
    \item Load the contract text from the configured file path.
    \item Run the full upstream pipeline (clause splitting, NER, SRL, intent
          classification).
    \item Embed all clause texts with \texttt{all-MiniLM-L6-v2}.
    \item Upsert embeddings, texts, and metadata into the ChromaDB collection.
    \item Display a success notification with the total number of indexed clauses.
\end{enumerate}

Once the index is built, query latency is dominated by embedding the question
(${\approx}10$\,ms) and the Gemini API round-trip, making the chatbot responsive
for interactive use.

\end{document}
